\chapter{Демонстрационный прототип и направления развития технологии}
\label{ch:prototype}

Настоящая глава связывает экспериментальную модель коллективных музыкальных рекомендаций, описанную в главе~\ref{ch:experiment}, с прикладным демонстрационным сценарием. Если главы~\ref{ch:sequential}--\ref{ch:experiment} формируют теоретическую и экспериментальную основу исследования, то данная глава показывает, каким образом полученный алгоритмический результат может быть представлен в демонстрационной сборке, приближенной к продуктовой концепции сервиса «Резонанс» / «Яндекс Музыка Резонанс».

\section{Назначение демонстрационного прототипа}
\label{sec:ch6:purpose}

В главе~\ref{ch:experiment} была разработана и оценена модель коллективных музыкальных рекомендаций для эфемерных групп пользователей. Экспериментальная постановка позволила проверить качество разных механизмов агрегации поверх фиксированного персонального скорера, однако сама по себе таблица метрик не показывает, как подобная модель может быть встроена в прикладной сценарий заведения или мероприятия. Поэтому после экспериментальной части целесообразно отдельно рассмотреть демонстрационный прототип, который иллюстрирует движение от исследовательского алгоритма к сервисной логике.

Демонстрационный прототип не является промышленной реализацией сервиса и не воспроизводит полный технологический контур «Яндекс Музыки». В нем не реализованы юридически значимые пользовательские согласия, B2B-лицензирование каталога, производственная инфраструктура публичного воспроизведения, интеграция с реальным профилем заведения и защищенный обмен данными внутри промышленной платформы. Его задача уже: показать принцип работы ключевой идеи, при которой состав присутствующей группы пользователей меняется, модель пересчитывает групповые рекомендации, а результат может быть обновлен в демонстрационном интерфейсе.

Важное ограничение связано с источником данных. В технической части работы используется открытый обезличенный датасет YAMBDA-50m. Идентификаторы пользователей и треков в нем анонимизированы; реальные названия композиций, исполнители, жанровые метки и человекочитаемые профили пользователей недоступны. Поэтому демонстрация не может показывать полноценный музыкальный плейлист в привычном потребительском смысле. Она оперирует идентификаторами пользователей и треков из датасета, а качество результата интерпретируется через численные метрики экспериментального протокола, прежде всего $\mathrm{NDCG@10}$ и $\mathrm{NDCG@20}$, а не через субъективное впечатление от конкретных известных композиций.

\section{Демонстрационный прототип модели коллективных музыкальных рекомендаций}
\label{sec:ch6:model-demo}

Техническая реализация проекта размещена в репозитории \texttt{music-recommendations-main}; внешняя версия репозитория указана в списке источников~\cite{music_recommendations_github}. На уровне алгоритмического слоя прототип использует ту же двухуровневую схему, что и экспериментальная глава: сначала персональный последовательный скорер формирует оценки треков для каждого пользователя, затем групповой агрегатор преобразует индивидуальные оценки в общий ранжированный список для эфемерной группы.

Фактический вход демонстрационного компонента задается как список идентификаторов пользователей из датасета YAMBDA. В коде этот контракт отражен в структуре \texttt{GroupSample}: поле \texttt{members} хранит кортеж raw \texttt{user\_id}, поле \texttt{candidates} содержит кандидатное множество треков, а поле \texttt{targets}~--- достижимые релевантные треки для расчета метрик. Кандидатное множество группы строится как объединение персональных top-200 кандидатов ее членов:
\begin{equation*}
    \mathcal{C}_G = \bigcup_{u \in G} \mathrm{Top}_{200}(u).
\end{equation*}
Значение 200 зафиксировано в конфигурации кэширования персональных скоров \texttt{InferenceConfig.K} и согласовано с экспериментальным протоколом главы~\ref{ch:experiment}. После формирования $\mathcal{C}_G$ агрегатор возвращает вектор групповых оценок по кандидатам; сортировка этих оценок дает ранжированный список анонимизированных \texttt{item\_idx}.

В локальной копии репозитория не обнаружен отдельный frontend-файл, в котором было бы жестко задано число строк, выводимых веб-интерфейсом как Top-N рекомендаций. Поэтому в настоящей работе не вводится произвольное значение такого параметра. По коду можно надежно зафиксировать другое: алгоритмический слой строит групповой ranking внутри кандидатного пула, сформированного из top-200 персональных кандидатов каждого участника, а стандартный модуль оценки возвращает метрики $\mathrm{NDCG@10}$ и $\mathrm{NDCG@20}$. Эти метрики реализованы в \texttt{src/eval/metrics.py} и \texttt{src/eval/group\_eval.py}, а их использование по умолчанию дополнительно зафиксировано в \texttt{GroupTrainConfig.eval\_k = (10, 20)}.

Следовательно, демонстрационный вывод может состоять из верхней части группового ранжирования и одной или нескольких метрик качества, рассчитанных в рамках того же протокола, что и в главе~\ref{ch:experiment}. При этом $\mathrm{NDCG@K}$ в такой демонстрации не является оценкой того, насколько конкретному человеку понравились анонимные треки. Метрика показывает, насколько построенный групповой ranking поднимает достижимые релевантные треки относительно тестовых прослушиваний членов синтетической группы. В силу обезличенности YAMBDA она выполняет роль количественного индикатора качества модели, а не полноценной пользовательской оценки музыкального опыта.

\section{Веб-демонстрация динамического изменения группы пользователей}
\label{sec:ch6:web-demo}

В продуктовой логике сервиса «Резонанс» музыкальный поток зависит не от абстрактной заранее заданной аудитории, а от фактического состава группы, присутствующей в заведении или на мероприятии. Поэтому демонстрационный сценарий должен показать динамику: пользователь появляется в условном пространстве, попадает в текущую группу, список \texttt{user\_id} обновляется, модель пересчитывает рекомендации, а интерфейс отображает новый результат.

Локальная копия репозитория подтверждает алгоритмический контракт такого сценария, но не содержит самостоятельного кода веб-интерфейса с таймером, визуализацией помещения или фиксированным интервалом обновления. Динамическое присутствие на уровне ML-слоя моделируется изменением списка пользователей, передаваемого в функцию построения \texttt{GroupSample}. Вспомогательный модуль \texttt{src/data/group\_synthesis.py} также содержит процедуру синтеза случайных эфемерных групп: из заданного пула пользователей выбираются группы размера от 2 до 5, без повторений внутри одной группы. Это соответствует исследовательской постановке, но не является подтверждением отдельного промышленного механизма появления и исчезновения реальных посетителей.

С учетом этих ограничений веб-демонстрацию корректно рассматривать как визуальную оболочку над алгоритмическим слоем. В таком сценарии на экране может отображаться условное пространство заведения или мероприятия, список текущих пользователей, блок групповых рекомендаций и блок метрики качества. Каждый отображаемый пользователь связан с некоторым идентификатором из YAMBDA. При изменении состава группы frontend должен передать новый список \texttt{user\_id} в рекомендательный компонент; далее на стороне модели заново формируются кандидаты, вычисляются групповые оценки и возвращается обновленный ranking анонимизированных треков.

Пересчет рекомендаций в таком сценарии должен запускаться по событию изменения состава группы или по явному запросу демонстрационного интерфейса. Конкретная частота автоматического обновления в локальном коде не зафиксирована, поэтому ее нельзя интерпретировать как экспериментально проверенную характеристику прототипа. Тем не менее сама возможность пересчета при изменении входного списка пользователей следует из структуры ML-компонентов: групповой агрегатор не привязан к постоянному идентификатору группы и получает на вход только текущих участников, их кандидатные множества и персональные скоры. Именно это свойство важно для концепции сервиса, поскольку оно позволяет рассматривать группу как эфемерную и изменяемую во времени.

\section{Демонстрационный прототип работы BLE-маяков}
\label{sec:ch6:ble-demo}

Вторая демонстрационная часть проекта относится не к ранжированию треков, а к прикладной идее фиксации присутствия пользователя в физическом пространстве. В папке \texttt{music-recommendations-main/BLE\_scanner\_demo} находится Android-приложение, демонстрирующее работу с BLE-маяками. Оно не связано напрямую с обучением модели и не передает данные в ML-пайплайн; его роль состоит в иллюстрации того, как в будущем может формироваться сигнал о принадлежности пользователя к зоне заведения или мероприятия.

Фактическая логика приложения подтверждается файлом \texttt{MainActivity.java} и README внутри папки прототипа. Пользователь может сохранить MAC-адрес BLE-маяка, назначить ему понятное имя и цвет. Список сохраненных маяков хранится локально в \texttt{SharedPreferences}. Приложение запрашивает разрешения, необходимые для Bluetooth-сканирования, запускает \texttt{BluetoothLeScanner} в режиме \texttt{SCAN\_MODE\_LOW\_LATENCY}, принимает результаты сканирования и обновляет для каждого найденного устройства значение RSSI и время последнего обнаружения.

Выбор ближайшего сохраненного маяка реализован сравнением свежих сигналов: приложение рассматривает только те сохраненные MAC-адреса, для которых сигнал был получен в течение последних 5 секунд, и выбирает маяк с максимальным RSSI. Более старые сигналы удаляются из карты наблюдений после 15 секунд. В интерфейсе показываются имя выбранного маяка, его MAC-адрес и текущее значение RSSI; верхняя панель окрашивается в цвет, заданный пользователем для этого маяка. Таким образом, прототип демонстрирует локальное распознавание ближайшего сохраненного BLE-маяка по относительной силе сигнала.

Принципиально важно, что такая логика не является системой точной геолокации. RSSI зависит от модели устройства, положения телефона, препятствий, отражений сигнала и плотности пространства, поэтому он может использоваться только как приближенный индикатор относительной близости к маяку. В контексте настоящей ВКР этот прототип показывает техническую реализуемость идеи «пользователь ближе к зоне А, чем к зоне Б», но не доказывает возможность точного трекинга координат посетителя и не заменяет отдельного исследования надежности BLE-инфраструктуры в реальном заведении.

\section{Направления дальнейшего развития технологии}
\label{sec:ch6:future}

Демонстрационная сборка фиксирует базовую связку между моделью группового ранжирования и прикладным сценарием, однако промышленная версия подобного сервиса потребовала бы существенно более сложной архитектуры. Ниже рассмотрены направления развития, которые связывают технический результат главы~\ref{ch:experiment} с продуктовой концепцией главы~\ref{ch:service-concept}.

\subsection{Фильтрация и бизнес-ограничения при формировании плейлиста}
\label{subsec:ch6:business-filters}

Групповые рекомендации не должны напрямую превращаться в публичный музыкальный эфир без учета контекста бизнеса. Заведение или организатор мероприятия должен иметь возможность задавать рамки музыкального потока: жанры, энергичность, формат площадки, желаемое настроение, стоп-листы треков или исполнителей, допустимый уровень популярности, ограничения по языку, тематике и иным параметрам. Иначе система может формально оптимизировать групповую релевантность, но нарушить брендовый стиль, правовой контур или сценарную драматургию пространства.

Простейший вариант реализации состоит в том, что коллективная модель работает не по всему каталогу, а внутри заранее разрешенного множества треков или редакторских каталогов. Более развитый вариант предполагает отдельный многоступенчатый слой кандидатной генерации, фильтрации и ранжирования. Такая логика согласуется с тем, как публично описывается персонализированный поток «Моей волны» в «Яндекс Музыке»: в статье Яндекса на Хабре показано, что поток строится через последовательность этапов, включая отбор кандидатов, фильтрацию и ранжирование, причем разные продуктовые режимы могут иметь собственные ограничения и приоритеты~\cite{stepurin2024unknown_my_wave}. Эта статья не описывает B2B-сервис или коллективные рекомендации, но служит примером того, что промышленная рекомендательная система музыкального сервиса обычно включает не только модель скоринга, но и отдельные продуктовые слои управления кандидатами.

На этой основе можно предположить, что промышленная версия сервиса коллективных музыкальных рекомендаций также должна содержать самостоятельный слой бизнес-ограничений. Он позволит совместить персонализацию под фактическую аудиторию с управляемостью музыкальной среды для площадки. В исследовательском прототипе такой слой не реализован: модель ранжирует кандидатов, сформированных из данных YAMBDA, и не знает о реальном жанре, тексте, исполнителе или коммерческом контексте трека.

\subsection{Взвешивание вклада отдельных пользователей}
\label{subsec:ch6:user-weights}

Базовая постановка групповой рекомендации предполагает агрегирование вклада всех участников группы. Однако в прикладном сервисе может потребоваться неравномерное влияние отдельных пользователей. Например, бизнес может доверять музыкальному вкусу определенного сотрудника, редактора или условного куратора, который фактически задает желаемую атмосферу заведения. В событийном сценарии повышенный вес может получать определенный сегмент посетителей, если именно его вовлеченность является важной продуктовой целью.

Такой механизм можно интерпретировать как управляемую модификацию группового профиля. С одной стороны, система сохраняет персонализацию под текущую аудиторию; с другой~--- не отказывается от авторского стиля и музыкальной политики площадки. В простейшей форме это может быть реализовано как веса пользователей в агрегирующей функции, а в более сложной~--- как отдельный контекстный слой, который меняет вклад участников в зависимости от зоны, времени, формата мероприятия или настроек бизнеса.

В текущей экспериментальной модели такой компонент не реализован и не проверялся. Поэтому взвешивание пользователей следует рассматривать как направление дальнейшего развития, а не как уже доказанную функциональность. Его внедрение потребует отдельной проверки: нужно будет оценить не только рост стандартных метрик ранжирования, но и влияние на справедливость по отношению к участникам группы.

\subsection{Несколько BLE-маяков и пространственное распределение аудитории}
\label{subsec:ch6:multi-ble}

Следующее направление связано с развитием BLE-сценария. В заведении или на мероприятии может быть установлено несколько маяков: у танцпола, у бара, в лобби, у сцены или в других функциональных зонах. Сравнение относительной близости пользователя к разным маякам может дать грубую оценку того, к какой зоне пространства он в данный момент ближе. Такая схема не должна трактоваться как точная геолокация; корректнее говорить о сценарной сегментации аудитории по зонам на основе приближенного сигнала.

Для мероприятий и вечеринок это открывает отдельную гипотезу. Посетители, находящиеся ближе к танцполу, могут быть условно более вовлечены в событийный сценарий; посетители ближе к бару или лобби могут находиться в менее активной зоне. Если система назначает больший вес пользователям из менее вовлеченной зоны, она теоретически может формировать рекомендации, направленные на их мягкую активацию: например, чаще выбирать знакомые или предпочтительные для этой группы треки, не разрушая общий стиль мероприятия.

В таком варианте сервис используется не только как генератор общего музыкального потока, но и как инструмент адаптации сценария мероприятия к пространственной структуре аудитории. Однако данная логика требует отдельной эмпирической проверки на реальных данных поведения посетителей. Необходимо учитывать шум RSSI, неоднозначность связи между зоной и вовлеченностью, юридические требования к обработке данных о присутствии, необходимость информированного согласия и приватностную архитектуру. В рамках текущей ВКР этот сценарий только концептуализируется и не проверяется экспериментально.

\subsection{Ассистент для диджея}
\label{subsec:ch6:dj-assistant}

Отдельное B2B-применение технологии связано с профессиональной работой DJ. Сервис может развиться не только как автоматический генератор музыкального потока, но и как аналитический ассистент. В таком сценарии DJ загружает подготовленный плейлист или набор треков, которые планирует использовать в сете, а система оценивает их потенциальную релевантность текущей аудитории мероприятия.

Если треки сопоставимы с каталогом платформы, оценка может происходить по каталожным идентификаторам. Если используются локальные аудиофайлы, перспективными становятся механизмы сопоставления через цифровые аудиоотпечатки, аудиоэмбеддинги или их комбинацию. Это позволило бы оценивать, насколько трек знаком аудитории, насколько он близок к ее музыкальным предпочтениям и какие блоки DJ-сета потенциально соответствуют текущему профилю гостей.

Такой ассистент не заменяет DJ и не подменяет профессиональную драматургию музыкального выступления. Его корректнее рассматривать как дополнительный аналитический сигнал, который помогает принять решение, но не автоматизирует художественный выбор полностью. В текущем эксперименте этот сценарий не реализован и не проверялся; он требует отдельной постановки задачи, отдельного интерфейса и проверки на данных реальных мероприятий.

\subsection{Новые данные о поведении аудитории и самостоятельный B2B-продукт аналитики}
\label{subsec:ch6:b2b-analytics}

Предлагаемый сервис потенциально создает новый тип данных на стыке музыкального потока, состава аудитории, изменений присутствующей группы, пространственного распределения посетителей и реакции бизнеса на разные музыкальные сценарии. При развитии сервиса можно исследовать взаимосвязь между параметрами музыкального потока, составом аудитории, перемещением посетителей по зонам пространства и поведением аудитории в заведении или на мероприятии.

В будущем подобные данные могут сопоставляться с агрегированными транзакционными или операционными показателями бизнеса, если это юридически допустимо, этически обосновано и методологически корректно. Тогда сервис может стать не только системой музыкальных рекомендаций, но и самостоятельной B2B-аналитической платформой для оценки реакции посетителей на аудиосреду. Ценность в таком случае имеют не только сами рекомендации, но и инструменты анализа того, как меняется аудитория и какие музыкальные сценарии связаны с разными режимами ее поведения.

Этот сценарий требует особенно аккуратной формулировки с точки зрения приватности. Нельзя исходить из того, что пользователи легко согласятся на непрозрачный мониторинг. Более корректная гипотеза состоит в другом: поскольку пользователь получает понятную ценность в виде адаптации музыкального опыта, получение информированного согласия на ограниченный и ясно описанный сценарий обработки данных может быть потенциально реалистичнее, чем для абстрактной системы наблюдения. Тем не менее подобное направление требует отдельной правовой проработки, этической оценки, аккуратного продуктового дизайна, эмпирической проверки спроса со стороны бизнеса и проверки готовности пользователей участвовать. Экономический потенциал аналитических модулей в рамках текущей ВКР не доказывается.

\section{Выводы по главе 6}
\label{sec:ch6:conclusions}

В главе показано, каким образом экспериментальная модель коллективных музыкальных рекомендаций может быть связана с прикладным демонстрационным сценарием сервиса. Алгоритмическая часть прототипа использует список идентификаторов пользователей YAMBDA, формирует эфемерную группу, строит кандидатное множество на основе top-200 персональных кандидатов каждого участника и возвращает ранжированный список анонимизированных треков. Метрики $\mathrm{NDCG@10}$ и $\mathrm{NDCG@20}$ интерпретируются как показатели качества группового ранжирования в рамках экспериментального протокола, а не как субъективная оценка музыкального опыта.

Отдельно рассмотрен Android BLE-прототип. Он демонстрирует сохранение BLE-маяков по MAC-адресам, сканирование окружающих Bluetooth-устройств и выбор ближайшего сохраненного маяка по самому сильному свежему RSSI-сигналу. Этот компонент не является системой точной геолокации и не реализует промышленный контур определения посетителей, но иллюстрирует техническую возможность локального распознавания относительной близости к BLE-маяку.

Дальнейшее развитие технологии связано с добавлением фильтров и бизнес-ограничений, взвешиванием вклада отдельных пользователей, использованием нескольких BLE-зон, разработкой ассистента для DJ и формированием B2B-аналитики на основе новых данных о музыкальном поведении аудитории. Эти направления расширяют исходную модель за пределы лабораторного эксперимента, но требуют самостоятельной эмпирической, правовой и продуктовой проверки.

% TODO(групповая ВКР): указать распределение задач между участниками проекта в рамках главы 6.
